%-------------------------
% Resume in Latex
% Author : Jake Gutierrez
% Based off of: https://github.com/sb2nov/resume
% License : MIT
% Content: Abdelrahman Khalifa; every fact is from ../../../facts.md.
% Generated by ../../../tools/make_jake_cv.py from ../Abdelrahman-Khalifa-CV.tex
%------------------------

\documentclass[letterpaper,11pt]{article}

\usepackage{latexsym}
\usepackage[empty]{fullpage}
\usepackage{titlesec}
\usepackage{marvosym}
\usepackage[usenames,dvipsnames]{color}
\usepackage{verbatim}
\usepackage{enumitem}
\usepackage[hidelinks]{hyperref}
\usepackage{fancyhdr}
\usepackage[english]{babel}
\usepackage{tabularx}
\input{glyphtounicode}
% Only change to the template's preamble: Latin Modern with T1 encoding. It looks
% the same as the default Computer Modern, but its bold font is named "...-Bold"
% (parsers such as OpenResume detect job titles by that) and "_" extracts as text.
\usepackage[T1]{fontenc}
\usepackage{lmodern}


%----------FONT OPTIONS----------
% sans-serif
% \usepackage[sfdefault]{FiraSans}
% \usepackage[sfdefault]{roboto}
% \usepackage[sfdefault]{noto-sans}
% \usepackage[default]{sourcesanspro}

% serif
% \usepackage{CormorantGaramond}
% \usepackage{charter}


\pagestyle{fancy}
\fancyhf{} % clear all header and footer fields
\fancyfoot{}
\renewcommand{\headrulewidth}{0pt}
\renewcommand{\footrulewidth}{0pt}

% Adjust margins
\addtolength{\oddsidemargin}{-0.5in}
\addtolength{\evensidemargin}{-0.5in}
\addtolength{\textwidth}{1in}
\addtolength{\topmargin}{-.5in}
\addtolength{\textheight}{1.0in}

\urlstyle{same}

\raggedbottom
\raggedright
\setlength{\tabcolsep}{0in}

% Sections formatting
\titleformat{\section}{
  \vspace{-4pt}\scshape\raggedright\large
}{}{0em}{}[\color{black}\titlerule \vspace{-5pt}]

% Ensure that generate pdf is machine readable/ATS parsable
\pdfgentounicode=1

%-------------------------
% Custom commands
\newcommand{\resumeItem}[1]{
  \item\small{
    {#1 \vspace{-2pt}}
  }
}

\newcommand{\resumeSubheading}[4]{
  \vspace{-2pt}\item
    \begin{tabular*}{0.97\textwidth}[t]{l@{\extracolsep{\fill}}r}
      \textbf{#1} & #2 \\
      \textit{\small#3} & \textit{\small #4} \\
    \end{tabular*}\vspace{-7pt}
}

\newcommand{\resumeSubSubheading}[2]{
    \item
    \begin{tabular*}{0.97\textwidth}{l@{\extracolsep{\fill}}r}
      \textit{\small#1} & \textit{\small #2} \\
    \end{tabular*}\vspace{-7pt}
}

\newcommand{\resumeProjectHeading}[2]{
    \item
    \begin{tabular*}{0.97\textwidth}{l@{\extracolsep{\fill}}r}
      \small#1 & #2 \\
    \end{tabular*}\vspace{-7pt}
}

\newcommand{\resumeSubItem}[1]{\resumeItem{#1}\vspace{-4pt}}

\renewcommand\labelitemii{$\vcenter{\hbox{\tiny$\bullet$}}$}

\newcommand{\resumeSubHeadingListStart}{\begin{itemize}[leftmargin=0.15in, label={}]}
\newcommand{\resumeSubHeadingListEnd}{\end{itemize}}
\newcommand{\resumeItemListStart}{\begin{itemize}}
\newcommand{\resumeItemListEnd}{\end{itemize}\vspace{-5pt}}

% Same link and code helpers as the main CVs
\newcommand{\cvlink}[2]{\href{#1}{\underline{#2}}}

%-------------------------------------------
%%%%%%  RESUME STARTS HERE  %%%%%%%%%%%%%%%%%%%%%%%%%%%%


\begin{document}

%----------HEADING----------
\begin{center}
    \textbf{\Huge \scshape Abdelrahman Khalifa} \\ \vspace{1pt}
    \small +20 114 250 5042 $|$ Cairo, Egypt $|$ \href{mailto:abdelrahman.khalifa.gad@gmail.com}{\underline{abdelrahman.khalifa.gad@gmail.com}} $|$
    \href{https://www.linkedin.com/in/abd002}{\underline{linkedin.com/in/abd002}} $|$
    \href{https://github.com/Abd002}{\underline{github.com/Abd002}}
\end{center}

%-----------EDUCATION-----------
\section{Education}
  \resumeSubHeadingListStart
    \resumeSubheading
      {Minia University}{Minia, Egypt}
      {Bachelor of Science (B.Sc.), Computer and Systems Engineering}{2020 -- 2025}
      \resumeItemListStart
        \resumeItem{Grade: Very Good with Honors. Graduation project: AUTOSAR V2X vehicle on STM32 with a dual-bank bootloader.}
        \resumeItem{Activities: Codeforces Expert; 4-time ECPC finalist (2021--2024); Treasurer, IEEE Student Branch.}
      \resumeItemListEnd
  \resumeSubHeadingListEnd

%-----------EXPERIENCE-----------
\section{Experience}
  \resumeSubHeadingListStart

    \resumeSubheading
      {Google Summer of Code Contributor}{May 2026 -- Oct 2026}
      {The Linux Foundation (OpenPrinting)}{Remote}
      \resumeItemListStart
        \resumeItem{Built the \textbf{Printers page for COSMIC Settings} and a standalone app: printers grouped by device, supply levels, print queues, live updates and an Add Printer flow for legacy printers via PAPPL (\cvlink{https://github.com/pop-os/cosmic-settings/pull/2182}{cosmic-settings \#2182}, in review).}
        \resumeItem{Designed the stack as 5 Rust crates: \texttt{printers-core} (types, grouping queues into physical devices), \texttt{printers-client} (async varlink client), \texttt{printers-server} (CUPS/IPP), \texttt{printers-ui} (iced) and \texttt{printers-app}.}
        \resumeItem{Ran the backend as a varlink IPC service in \texttt{cosmic-settings-daemon} (\cvlink{https://github.com/pop-os/cosmic-settings-daemon/pull/188}{cosmic-settings-daemon \#188}, in review), so the C FFI code never loads in the Settings process; the standalone app falls back to running it in-process.}
      \resumeItemListEnd

    \resumeSubheading
      {Embedded Software Engineer}{Aug 2025 -- Present}
      {EKSON}{Egypt}
      \resumeItemListStart
        \resumeItem{Wrote the ESP32-S3 firmware and migrated it from ESP-IDF/FreeRTOS to Zephyr RTOS: PID valve control with IMU feedback, safety stops and a UART protocol.}
        \resumeItem{Developed the PyQt5 desktop control app, cutting CPU usage and latency by about 40\% (measured) using asyncio to wait for ESP32 and server messages concurrently; added remote control and firmware updates through a Node.js/MongoDB backend.}
      \resumeItemListEnd

    \resumeSubheading
      {Embedded Linux Systems Intern}{Dec 2024 -- Jan 2025}
      {Orange Digital Center (ODC)}{Egypt}
      \resumeItemListStart
        \resumeItem{Created custom Linux images with Yocto and Buildroot and wrote a Linux kernel module; worked with device trees, root file systems and hardware flashing. Scored 100\% in the final evaluation.}
      \resumeItemListEnd

  \resumeSubHeadingListEnd

%-----------PROJECTS-----------
\section{Projects}
    \resumeSubHeadingListStart
      \resumeProjectHeading
          {\textbf{OpenPrinting cups-rs (Rust) and libcups (C)} $|$ \emph{Upstream contributions (all accepted)}}{2026}
          \resumeItemListStart
            \resumeItem{cups-rs: ported to libcups3 (CUPS 2 and 3 builds, CI); added IPP management operations, Printer Application support and FFI bindings for the libcups3 DNS-SD API (\cvlink{https://github.com/OpenPrinting/cups-rs/pull/18}{PRs \#18--\#22}).}
            \resumeItem{libcups: fixed crashes in \texttt{cupsDNSSDDelete()} and a leaked client in the Avahi DNS-SD code (\cvlink{https://github.com/OpenPrinting/libcups/pull/164}{\#164}, \cvlink{https://github.com/OpenPrinting/libcups/pull/165}{\#165}).}
          \resumeItemListEnd
      \resumeProjectHeading
          {\textbf{Zephyr RTOS, Boa JavaScript engine (Rust), Rust Clippy} $|$ \emph{Contributor (all merged)}}{Mar 2025 -- Dec 2025}
          \resumeItemListStart
            \resumeItem{Zephyr RTOS: added the \texttt{IP\_MULTICAST\_LOOP} socket option for IPv4 in 6 networking modules (\cvlink{https://github.com/zephyrproject-rtos/zephyr/issues/86977}{\#86977}).}
            \resumeItem{Boa: optimized \texttt{JsStr::code\_points()} for Latin1 strings, avoiding UTF-16 conversion (\cvlink{https://github.com/boa-dev/boa/pull/4553}{PR \#4553}).}
            \resumeItem{Rust Clippy: fixed \texttt{ref\_as\_ptr} false positives for non-temporary references in \texttt{let}/\texttt{const} (\cvlink{https://github.com/rust-lang/rust-clippy/pull/16214}{PR \#16214}).}
          \resumeItemListEnd
      \resumeProjectHeading
          {\textbf{Kernel Monitor} $|$ \emph{C, Linux, ARM, QEMU}}{\cvlink{https://github.com/Abd002/kernel-monitor}{GitHub}}
          \resumeItemListStart
            \resumeItem{Character device driver that walks the process list under RCU and exports per-task IDs, memory, threads, context switches and CPU time to user space; run on QEMU.}
          \resumeItemListEnd
      \resumeProjectHeading
          {\textbf{IoT Device Communication System} $|$ \emph{C++, Yocto}}{\cvlink{https://github.com/Abd002/IoT-Comm-System-RPi4-QEMU}{GitHub}}
          \resumeItemListStart
            \resumeItem{TCP and UDP multicast between a Pi and QEMU clients.}
          \resumeItemListEnd
    \resumeSubHeadingListEnd

%
%-----------PROGRAMMING SKILLS-----------
\section{Technical Skills}
 \begin{itemize}[leftmargin=0.15in, label={}]
    \small{\item{
     \textbf{Languages}{: Rust, C, Modern C++, Python (asyncio), Bash, JavaScript (Node.js)} \\
     \textbf{Systems \& Linux}{: Linux (Ubuntu/Debian), operating systems, FFI and bindings, IPC (varlink, D-Bus; zbus, zlink), concurrency (async/event-driven, multithreading), memory management, kernel modules and device drivers} \\
     \textbf{Networking}{: TCP/UDP sockets, IPv4 multicast, HTTP/REST, IPP, DNS-SD/mDNS (Avahi), CUPS} \\
     \textbf{Tooling}{: Cargo, Clippy, rustfmt, Git/GitHub, GitHub Actions CI, CMake, Make, Yocto, Buildroot, QEMU, testing, debugging, performance optimization, code review}
    }}
 \end{itemize}


%-------------------------------------------
\end{document}
